\begin{figure}[!htbp]
\centering
\begin{tikzpicture}[font=\small\sffamily,x=0.85cm,y=0.85cm]
\foreach \shift/\heading in {0/{8-path: ambiguous},5/{m-path: source illustration}}{
\begin{scope}[xshift=\shift cm]
\foreach \x/\y in {0/2,1/2,0/1,1/0}{\draw[MidGrey,fill=PaleRed] (\x,\y) rectangle ++(1,1);}
\node[above] at (1,3.4) {\heading};
\node[right] at (2,2.5) {$p$}; \node[below] at (1.5,0) {$q$};
\draw[cvarrow] (1.5,2.5)--(0.5,2.5)--(0.5,1.5)--(1.5,0.5);
\end{scope}}
\draw[SapienzaRed,dashed,line width=0.9pt] (1.5,2.5)--(0.5,1.5);
\end{tikzpicture}
\caption{The source example removes the alternative diagonal route in its m-path illustration, retaining a single route from $p$ to $q$.}
\label{fig:connectivity-paths}
\end{figure}
\begin{remarkbox}{Source rule and illustration}
The displayed rule uses $N_4(p)\cap N_4(q)=\varnothing$ without restricting the intersection to the selected pixel class. On the full interior grid, diagonal pixels share two 4-neighbors, so the literal rule excludes all such diagonal links. The source's path illustration nevertheless retains the final diagonal step. Both are preserved here: the diagram reproduces the slide's intended ambiguity removal, rather than silently replacing its written definition with another convention.
\end{remarkbox}
